\providecommand{\EinDepth}[1]{\operatorname{Ein}_{#1}}
\providecommand{\Eexp}{\mathcal E_{\mathrm{exp}}}

\section{Multipoint one-level polyexponentials of arbitrary weight}
\label{sec:ein-depth-tower}

For \(r\geq1\), define the order-\(r\) complementary exponential
integral by
\begin{equation}\label{eq:ein-depth-definition}
 \EinDepth{r}(z)
 :=\sum_{n\geq1}\frac{(-1)^{n-1}z^n}{n^r n!}.
\end{equation}
We call \(r\) the \emph{order} (or weight), not the level: these are
classical one-level polyexponentials, not the multiple-index
polyexponentials of level greater than one
\cite{Hardy1905PartII,Boyadzhiev2007,AminovArnaudo2026}.
Thus \(\EinDepth{1}=\Ein\), and
\begin{equation}\label{eq:ein-depth-chain}
 z\EinDepth{1}'(z)=1-e^{-z},\qquad
 z\EinDepth{r}'(z)=\EinDepth{r-1}(z)\quad(r\geq2).
\end{equation}
We use the harmless notation
\(\EinDepth{0}(z)=1-e^{-z}\)
only when writing the chain compactly.  For each fixed \(r\),
\(\EinDepth{r}\) is an \(E\)-function.  Indeed, its coefficient sequence
is hypergeometric, hence holonomic.  Moreover, for every \(N\), the single
integer \(\operatorname{lcm}(1,\ldots,N)^r\) clears simultaneously all
coefficients \(n![z^n]\EinDepth{r}(z)\), \(0\leq n\leq N\); it grows at
most exponentially in \(N\).  The coefficients are rational of absolute
value at most one, so the conjugate-growth condition is automatic.

We first record the rational-denominator step needed in the iterated
primitive argument.  It is useful beyond the present family.

\begin{lemma}[Affine form of a primitive rational function]
\label{lem:depth-affine-primitive}
Let \(F\) be a differential field of characteristic zero whose constant
field \(C\) is algebraically closed.  Let
\(a_1,\ldots,a_s\in F\) have the property
\begin{equation}\label{eq:depth-no-base-exact-combination}
 \sum_{i=1}^s c_i a_i\in D(F),\quad c_i\in C
 \qquad\Longrightarrow\qquad c_1=\cdots=c_s=0.
\end{equation}
Adjoin algebraically independent primitives
\(u_1,\ldots,u_s\), with \(Du_i=a_i\), and put
\(L=F(u_1,\ldots,u_s)\).  If \(g\in L\) and \(Dg\in F\), then
\begin{equation}\label{eq:depth-affine-conclusion}
 g=f+\sum_{i=1}^s b_i u_i
 \qquad(f\in F,\ b_i\in C).
\end{equation}
In particular, the constant field of \(L\) is again \(C\).
\end{lemma}

\begin{proof}
Work first in \(F[u_1,\ldots,u_s]\).  There is no nonconstant
irreducible Darboux polynomial in this ring.  Indeed, suppose that an
irreducible nonconstant \(P\) divides \(DP\).  Degree in the \(u_i\)
shows that \(DP=hP\) for some \(h\in F\).  Fix a monomial order refining
total degree and divide \(P\) by the coefficient of its leading monomial.
For the resulting monic polynomial \(Q\), one has
\(DQ=(h-Dp/p)Q\), where \(p\) was that leading coefficient.  Since the
derivation lowers the \(u\)-degree of the monomial part and the leading
coefficient of \(Q\) is one, comparison of the leading monomial forces
\(h-Dp/p=0\); hence \(DQ=0\).
Write \(Q_n\) for its highest nonzero homogeneous part.  The degree-\(n\)
part of \(DQ=0\) shows that every coefficient of \(Q_n\) belongs to \(C\).
The degree-\((n-1)\) part is
\[
 DQ_{n-1}+\sum_{i=1}^s
 a_i\frac{\partial Q_n}{\partial u_i}=0.
\]
Since \(Q_n\) is nonconstant, at least one coefficient comparison in this
identity writes a nonzero constant linear combination of the \(a_i\) as
an element of \(D(F)\), contrary to
\eqref{eq:depth-no-base-exact-combination}.

Now write \(g=A/B\) in lowest terms.  If a nonconstant irreducible factor
\(P\) occurs in \(B\), put \(m=-v_P(g)>0\).  If \(P\nmid DP\), the usual
quotient-rule computation in the discrete valuation \(v_P\) gives
\(v_P(Dg)=-m-1\): the term obtained by differentiating \(P^{-m}\) has
that valuation and cannot cancel with terms of valuation at least \(-m\).
But every nonzero element of \(F\) has \(P\)-valuation zero.  Thus
\(Dg\in F\) forces \(P\mid DP\), which the preceding paragraph rules out.
Hence \(g\) is a polynomial in the \(u_i\).

Let \(g_n\) be its highest homogeneous part.  Since \(Dg\in F\), the
degree-\(n\) coefficient comparison makes the coefficients of \(g_n\)
constant.  If \(n\geq2\), the degree-\((n-1)\) comparison again makes a
nonzero constant linear combination of the \(a_i\) exact in \(F\), a
contradiction.  Thus \(n\leq1\), and
\eqref{eq:depth-affine-conclusion} follows.  If \(Dg=0\), then
\(Df+\sum_i b_i a_i=0\);
\eqref{eq:depth-no-base-exact-combination} gives \(b_i=0\), and then
\(f\in C\).  Thus \(L\) has no new constants.
\end{proof}

We now apply this lemma to the exponential torus used in
Lemma~\ref{lem:exact}.  Let
\(\lambda_1,\ldots,\lambda_s\in\Qbar^\times\) be distinct, let
\(\nu_1,\ldots,\nu_d\) be a \(\mathbb Z\)-basis of the additive group
generated by the \(\lambda_i\), and write
\(\lambda_i=m_i\mathbin{\cdot}\nu\), with
\(m_i\in\mathbb Z^d\).  Put
\[
 K=\operatorname{Frac}
 \bigl(\Qbar(z)[X_1^{\pm1},\ldots,X_d^{\pm1}]\bigr),
 \qquad DX_j=-\nu_jX_j.
\]
Because the \(\nu_j\) form a basis of the group generated by the
\(\lambda_i\), each \(\nu_j\) is also an integral combination of the
\(\lambda_i\).  Thus passing from \(\lambda\) to \(\nu\) changes neither
the generated exponential field nor its character lattice.  The vectors
\(m_i\) are pairwise distinct and nonzero.
The constant field of \(K\) is \(\Qbar\).  Lemma~\ref{lem:exact}
implies both
\begin{equation}\label{eq:depth-base-nonexact}
 \frac1z\notin D(K)
\end{equation}
and
\begin{equation}\label{eq:depth-first-layer-no-combination}
 \sum_{i=1}^s c_i\frac{1-X^{m_i}}z\in D(K),\quad c_i\in\Qbar
 \qquad\Longrightarrow\qquad c_1=\cdots=c_s=0.
\end{equation}
Indeed, after collecting the distinct Laurent characters, each assertion
is an immediate instance of that lemma.

All the functions below have Taylor coefficients in \(\Qbar\), and the
field they generate at the origin embeds in \(\Qbar((z))\).  Since the
constant field of \(\Qbar((z))\) under \(d/dz\) is \(\Qbar\), the ambient
extension used in every invocation of Kolchin--Ostrowski has no new
constants.  The affine lemma below also verifies this internally after each
primitive layer is adjoined.

For \(q\geq1\), set
\[
 Y_{i,q}(z)=\EinDepth{q}(\lambda_i z),
 \qquad
 K_q=K(Y_{i,j}:1\leq i\leq s,\ 1\leq j\leq q).
\]
Then
\begin{equation}\label{eq:depth-tower-derivatives}
 DY_{i,1}=\frac{1-X^{m_i}}z,
 \qquad
 DY_{i,q}=\frac{Y_{i,q-1}}z\quad(q\geq2).
\end{equation}

\begin{lemma}[Nonexactness throughout the tower]
\label{lem:depth-one-over-z-nonexact}
For every \(q\geq0\), with \(K_0=K\), the constant field of \(K_q\) is
\(\Qbar\), and
\begin{equation}\label{eq:depth-one-over-z-all-levels}
 \frac1z\notin D(K_q).
\end{equation}
Moreover, at each level the variables
\(Y_{1,q},\ldots,Y_{s,q}\) are algebraically independent over
\(K_{q-1}\).
\end{lemma}

\begin{proof}
The primitive form of the Kolchin--Ostrowski theorem~\cite{Kolchin} says that primitives
with derivatives in a differential field with algebraically closed
constants are algebraically dependent only if a nonzero constant linear
combination of their derivatives is exact in the base field.
Equation~\eqref{eq:depth-first-layer-no-combination} therefore makes the
first layer algebraically independent over \(K\).  Applying
Lemma~\ref{lem:depth-affine-primitive} also shows that \(K_1\) has no new
constants.

Suppose \(Dg=1/z\) with \(g\in K_1\).  Since \(Dg\in K\),
Lemma~\ref{lem:depth-affine-primitive} gives
\(g=k+\sum_i b_iY_{i,1}\), where \(k\in K\) and
\(b_i\in\Qbar\).  Hence
\[
 Dk=\frac{1-\sum_i b_i+\sum_i b_iX^{m_i}}z,
\]
contradicting Lemma~\ref{lem:exact}.  This proves
\eqref{eq:depth-one-over-z-all-levels} for \(q=1\).

Assume the assertions through level \(q\geq1\).  The extension
\(K_q/K_{q-1}\) admits all translations
\[
 \tau_h(Y_{i,q})=Y_{i,q}+h_i,\qquad
 h=(h_1,\ldots,h_s)\in\Qbar^s,
\]
as differential automorphisms.  If, for some \(g\in K_q\),
\[
 Dg=\sum_{i=1}^s c_i\frac{Y_{i,q}}z
\]
with \(c\ne0\), choose \(h\) with \(c\mathbin{\cdot}h\ne0\).  Then
\[
 D(\tau_hg-g)=\frac{c\mathbin{\cdot}h}{z},
\]
contradicting the inductive nonexactness of \(1/z\) in \(K_q\).
Thus no nonzero constant linear combination of the derivatives of the
next layer is exact.  Kolchin--Ostrowski gives the algebraic independence
of \(Y_{1,q+1},\ldots,Y_{s,q+1}\) over \(K_q\), and
Lemma~\ref{lem:depth-affine-primitive} shows that the constants remain
\(\Qbar\).

Finally, suppose \(Dg=1/z\) with \(g\in K_{q+1}\).  The affine lemma gives
\[
 g=k+\sum_i b_iY_{i,q+1},\qquad
 k\in K_q,\quad b_i\in\Qbar,
\]
and therefore
\[
 Dk=\frac{1-\sum_i b_iY_{i,q}}z.
\]
If \(b\ne0\), translate the \(q\)-th layer by an \(h\) with
\(b\mathbin{\cdot}h\ne0\); the difference makes a nonzero multiple of
\(1/z\) exact in \(K_q\), a contradiction.  If \(b=0\), the displayed
identity directly contradicts the inductive hypothesis.  This completes
the induction.
\end{proof}

For clarity, the first new case can also be checked without the induction.
For one nonzero algebraic \(\lambda\), put
\(Y=\Ein(\lambda z)\) and
\(Z=\EinDepth{2}(\lambda z)\).  Exactness in \(K\) first shows that
\(Y\) is transcendental over \(K\); the affine lemma then gives constants
\(\Qbar\) in \(K(Y)\), and its displayed calculation proves
\(1/z\notin D(K(Y))\).  If \(Z\) were algebraic over \(K(Y)\), the
one-primitive Kolchin--Ostrowski theorem would give
\(Dg=Y/z\) for some \(g\in K(Y)\), after rescaling a nonzero constant.
Because \(Y\) is transcendental, \(\tau_h(Y)=Y+h\) is a differential
automorphism of \(K(Y)/K\).  Applying it with \(h\ne0\) gives
\(D(\tau_hg-g)=h/z\), the required contradiction.  Thus
\(e^{\lambda z},Y,Z\) have functional transcendence degree three when
\(s=1\), and the order-two claim is independent of any higher-order
induction.

\begin{theorem}[Multipoint resonant polyexponential theorem]
\label{thm:ein-depth-multipoint}
Let \(\lambda_1,\ldots,\lambda_s\in\Qbar^\times\) be distinct, let
\[
 d=\dim_{\Q}\operatorname{span}_{\Q}
   \{\lambda_1,\ldots,\lambda_s\},
\]
and let \(R\geq1\).  Then
\begin{equation}\label{eq:ein-depth-value-trdeg}
 \operatorname{trdeg}_{\Qbar}
 \Qbar\bigl(
   e^{\lambda_i},\EinDepth{r}(\lambda_i):
   1\leq i\leq s,\ 1\leq r\leq R
 \bigr)
 =d+Rs.
\end{equation}
In particular, the \(Rs\) polyexponential values are algebraically independent over
\(\Qbar(e^{\lambda_1},\ldots,e^{\lambda_s})\).  All algebraic relations
among the displayed exponential and polyexponential coordinates are therefore the
toric relations among the exponential coordinates.
\end{theorem}

\begin{proof}
Lemma~\ref{lem:depth-one-over-z-nonexact} gives functional transcendence
degree \(d+Rs\) over \(\Qbar(z)\): the exponential torus contributes
\(d\), and each of the \(R\) primitive layers contributes \(s\).

The same functional transcendence degree holds over \(\C(z)\).  Indeed,
all Taylor coefficients are algebraic; after clearing rational-function
denominators, any finite relation over \(\C(z)\) would be a finite linear
dependence over \(\C\) among series with coefficients in \(\Qbar\), and
rank invariance under the scalar extension \(\Qbar\subset\C\) would produce
a relation over \(\Qbar(z)\).

For specialization, include in one vector the constant function \(1\),
the finitely many exponentials \(e^{\pm\lambda_i z}\), and all
\(\EinDepth{r}(\lambda_i z)\), \(1\leq r\leq R\).  The exponentials are
\(E\)-functions.  If \(q_i\in\mathbb Z_{>0}\) clears the algebraic
denominator of \(\lambda_i\), then for every \(N\) the integer
\[
 q_i^N\operatorname{lcm}(1,\ldots,N)^r
\]
clears simultaneously the first \(N\) normalized coefficients of
\(\EinDepth{r}(\lambda_i z)\); all of their conjugates have size at most
\(C_i^N\) for a fixed \(C_i\).  Thus these scaled order functions satisfy
the denominator and conjugate-growth axioms for \(E\)-functions.  By
\eqref{eq:ein-depth-chain}, the displayed vector satisfies a first-order
linear system over \(\Qbar(z)\), with diagonal constant entries for the
exponentials and only \(1/z\) entries in the order chains.  Consequently
its coefficient matrix has no finite pole except possibly \(z=0\), so the
nonzero algebraic point \(z=1\) is ordinary.  Beukers'
refined Siegel--Shidlovskii theorem~\cite{Beukers2006} transfers the functional
transcendence degree, and all relations, to the values at \(z=1\).  This
proves \eqref{eq:ein-depth-value-trdeg}; the final statement follows from
the algebraic independence of all polyexponential coordinates over the exponential
field.
\end{proof}

\begin{corollary}[Pure exponential base change in every order]
\label{cor:ein-depth-pure-exponential}
The countable family
\[
 \bigl\{\EinDepth{r}(\lambda):
      \lambda\in\Qbar^\times,\ r\geq1\bigr\}
\]
is algebraically independent over
\(\Eexp=\Qbar(e^\alpha:\alpha\in\Qbar)\), in the finite-subset sense.
\end{corollary}

\begin{proof}
Any proposed polynomial relation uses finitely many polyexponential values and
finitely many coefficients from \(\Eexp\).  Enlarge the finite set of
algebraic points to include the exponents occurring in those coefficients,
take the largest order appearing, clear coefficient denominators, and
apply Theorem~\ref{thm:ein-depth-multipoint}.  The selected order values
remain algebraically independent over the enlarged exponential field.
\end{proof}

\begin{corollary}[Hausdorff moments and consecutive finite spectra]
\label{cor:ein-depth-hausdorff-spectra}
Fix \(a,b\in\Qbar\cap\mathbb R_{>0}\) and an integer \(\ell\geq1\), and
write \(\Delta_a f(s)=f(s+a)-f(s)\).  For \(q\geq1\) and \(m\geq0\), put
\begin{equation}\label{eq:ein-depth-finite-difference-moments}
 \mu_m^{(q)}
 :=(-1)^{\ell-1}(\Delta_a^\ell\EinDepth{q})(b+am).
\end{equation}
Then \((\mu_m^{(q)})_{m\geq0}\) is a strict Hausdorff moment sequence on
\([e^{-a},1]\), with
\begin{equation}\label{eq:ein-depth-hausdorff-density}
 \mu_m^{(q)}=\int_{e^{-a}}^1x^m w_q(x)\,dx,
 \qquad
 w_q(x)=
 \frac{x^{b/a-1}(1-x)^\ell}{(q-1)!(-\log x)}
 \left(\log\frac{a}{-\log x}\right)^{q-1}.
\end{equation}
The entire countable family
\begin{equation}\label{eq:ein-depth-all-moments-independent}
 \{\mu_m^{(q)}:q\geq1,\ m\geq0\}
\end{equation}
is algebraically independent over \(\Eexp\).

For each \(q\), let
\[
 \mathcal H_n^{(q)}
  =\det(\mu_{i+j}^{(q)})_{0\leq i,j\leq n},
 \qquad \mathcal H_{-1}^{(q)}=1.
\]
All these Hankel determinants, jointly for \(q\geq1\) and \(n\geq0\),
are algebraically independent over \(\Eexp\), and they are strictly
positive.  Let the monic orthogonal polynomials for \(w_q(x)\,dx\) satisfy
\begin{equation}\label{eq:ein-depth-favard-recurrence}
 P_{n+1}^{(q)}(x)
 =(x-\alpha_n^{(q)})P_n^{(q)}(x)
  -\beta_n^{(q)}P_{n-1}^{(q)}(x),
 \qquad \beta_n^{(q)}>0\quad(n\geq1).
\end{equation}
Then the family
\begin{equation}\label{eq:ein-depth-favard-independent}
 \{\mu_0^{(q)},\alpha_n^{(q)},\beta_{n+1}^{(q)}:
        q\geq1,\ n\geq0\}
\end{equation}
is algebraically independent over \(\Eexp\).

For every fixed \(q\), the infinite Hankel matrix
\((\mu_{i+j}^{(q)})_{i,j\geq0}\) is strictly totally positive.  The
bounded Jacobi operator associated with
\eqref{eq:ein-depth-favard-recurrence} has simple, purely absolutely
continuous spectrum \([e^{-a},1]\) of multiplicity one, and
\begin{equation}\label{eq:ein-depth-nevai-limits}
 \alpha_n^{(q)}\longrightarrow\frac{1+e^{-a}}2,
 \qquad
 \beta_n^{(q)}\longrightarrow\frac{(1-e^{-a})^2}{16}.
\end{equation}
Consequently
\begin{equation}\label{eq:ein-depth-hankel-capacity}
 \lim_{n\to\infty}(\mathcal H_n^{(q)})^{1/n^2}
 =\frac{1-e^{-a}}4,
\end{equation}
and the normalized zero-counting measures of the monic orthogonal
polynomials \(P_N^{(q)}\) converge weakly to
\begin{equation}\label{eq:ein-depth-arcsine-law}
 \frac{\mathbf 1_{(e^{-a},1)}(x)}
 {\pi\sqrt{(x-e^{-a})(1-x)}}\,dx.
\end{equation}

Finally, let \(J_N^{(q)}\) be the \(N\)-by-\(N\) Jacobi truncation attached
to \eqref{eq:ein-depth-favard-recurrence}.  For every fixed \(q\) and
\(N\geq2\), the \(2N-1\) eigenvalues of
\(J_N^{(q)}\) and \(J_{N-1}^{(q)}\), taken together, are algebraically
independent over \(\Eexp\).  The same joint assertion holds for finitely
many distinct orders \(q\), with one chosen consecutive pair at each
order.  Each spectrum is simple and the two consecutive spectra strictly
interlace.  The count \(2N-1\) is sharp: if \(\Sigma_j^{(q)}\) denotes
the full spectrum of \(J_j^{(q)}\), then
\begin{equation}\label{eq:ein-depth-full-finite-spectral-rank}
 \operatorname{trdeg}_{\Eexp}
 \Eexp\bigl(\Sigma_1^{(q)},\ldots,\Sigma_N^{(q)}\bigr)=2N-1.
\end{equation}
Thus the consecutive pair
\(\Sigma_{N-1}^{(q)}\cup\Sigma_N^{(q)}\) is a transcendence basis for
the full finite spectral tower up to algebraic extension.
\end{corollary}

\begin{proof}
The Mellin identity
\[
 \frac1{n^q}=\frac1{(q-1)!}
 \int_0^1t^{n-1}(-\log t)^{q-1}\,dt
\]
and termwise integration give, first near the origin and then everywhere
by analyticity,
\begin{equation}\label{eq:ein-depth-mellin-integral}
 \EinDepth{q}(s)=\frac1{(q-1)!}
 \int_0^1(1-e^{-st})(-\log t)^{q-1}\frac{dt}{t}.
\end{equation}
Since
\[
 (-1)^{\ell-1}\Delta_a^\ell(1-e^{-st})
 =e^{-st}(1-e^{-at})^\ell,
\]
substitution of \(s=b+am\), followed by \(x=e^{-at}\), gives exactly
\eqref{eq:ein-depth-hausdorff-density}.  Its density is positive on
\((e^{-a},1)\).  At \(x=e^{-a}\) the logarithmic factor vanishes when
\(q>1\) (and is interpreted as one when \(q=1\)); as \(x\uparrow1\),
\[
 w_q(x)=O\!\left((1-x)^{\ell-1}
             \left|\log(1-x)\right|^{q-1}\right).
\]
Thus the density is integrable at both endpoints and has infinite support.
This proves the strict Hausdorff assertion.

For fixed \(q\), every finite set of the \(\mu_m^{(q)}\) is obtained from
the values
\(\EinDepth{q}(b+ak)\) by a matrix whose row indexed by \(m\) has last
nonzero entry, with nonzero coefficient, in column \(m+\ell\).  Distinct
rows therefore have distinct pivots.  For several orders the matrix is
block diagonal in \(q\).  Corollary~\ref{cor:ein-depth-pure-exponential}
then proves \eqref{eq:ein-depth-all-moments-independent}.

Strict positivity of the measure gives \(\mathcal H_n^{(q)}>0\).  Moreover,
\begin{equation}\label{eq:ein-depth-hankel-triangular-jacobian}
 \frac{\partial\mathcal H_n^{(q)}}
      {\partial\mu_{2n}^{(q)}}=\mathcal H_{n-1}^{(q)}.
\end{equation}
The Jacobian with respect to
\(\mu_0^{(q)},\mu_2^{(q)},\ldots,\mu_{2n}^{(q)}\) is triangular with
nonzero diagonal; together with the independence of the moments, this
proves the assertion about all Hankel determinants.

The finite Stieltjes--Favard algorithm gives, for every \(N\geq1\), the
birational equality
\begin{multline}\label{eq:ein-depth-moment-favard-fields}
 \Eexp(\mu_0^{(q)},\ldots,\mu_{2N-1}^{(q)})\\
 =\Eexp(\mu_0^{(q)},
   \alpha_0^{(q)},\ldots,\alpha_{N-1}^{(q)},
   \beta_1^{(q)},\ldots,\beta_{N-1}^{(q)}).
\end{multline}
Both sides have \(2N\) displayed generators.  Hence
\eqref{eq:ein-depth-favard-independent} follows, jointly in all orders.

Strict total positivity follows from Andreief's identity: every minor is
an integral of the product of two generalized Vandermonde determinants
against the positive density in
\eqref{eq:ein-depth-hausdorff-density}.  After ordering the integration
variables, both determinants have the same strict sign.  Normalizing the
measure by its mass identifies the infinite Jacobi operator with
multiplication by \(x\) on its \(L^2\)-space.  Since \(w_q\) is positive
almost everywhere on the whole interval, this gives simple purely
absolutely continuous spectrum \([e^{-a},1]\).  The
Denisov--Rakhmanov theorem~\cite{Denisov2004}, affinely rescaled to that
interval, gives \eqref{eq:ein-depth-nevai-limits}.  The identity
\[
 \mathcal H_n^{(q)}
 =(\mu_0^{(q)})^{n+1}
   \prod_{k=1}^n(\beta_k^{(q)})^{n+1-k}
\]
then proves \eqref{eq:ein-depth-hankel-capacity}; the standard fixed-walk
trace argument for Jacobi truncations with convergent coefficients gives
the arcsine law \eqref{eq:ein-depth-arcsine-law}.

The characteristic polynomial of \(J_N^{(q)}\) is \(P_N^{(q)}\).  Given
the two monic polynomials \(P_N^{(q)}\) and \(P_{N-1}^{(q)}\), successive
Euclidean divisions in \eqref{eq:ein-depth-favard-recurrence} recover
\(\alpha_0^{(q)},\ldots,\alpha_{N-1}^{(q)}\) and
\(\beta_1^{(q)},\ldots,\beta_{N-1}^{(q)}\) rationally; conversely the
recurrence constructs the two polynomials.  Their \(2N-1\) nonleading
coefficients are therefore algebraically independent.  Adjoining all
roots is an algebraic extension and introduces exactly \(2N-1\) displayed
roots, so those roots are algebraically independent as claimed.  Positivity
of the \(\beta_n^{(q)}\) gives simplicity and strict interlacing.
The same Euclidean divisions recover every earlier monic polynomial
\(P_j^{(q)}\), \(j<N-1\), from the coefficients of
\(P_N^{(q)}\) and \(P_{N-1}^{(q)}\).  All earlier eigenvalues are
therefore algebraic over the field generated by the consecutive pair,
which proves \eqref{eq:ein-depth-full-finite-spectral-rank}.
\end{proof}

\begin{corollary}[The order-two Basel split]
\label{cor:ein-depth-basel}
Put
\[
 \eta=\Ein(1)=\gamma+\frac{\delta}{e},\qquad
 T=\int_1^\infty\frac{e^{-u}\log u}{u}\,du.
\]
Then each of the following triples is algebraically independent over
\(\Qbar\):
\begin{equation}\label{eq:ein-depth-basel-independent-triples}
 \bigl(e,\eta,\EinDepth{2}(1)\bigr),
 \qquad
 \left(e,\gamma+\frac{\delta}{e},
       \gamma^2+\frac{\pi^2}{6}-2T\right).
\end{equation}
Consequently
\begin{equation}\label{eq:ein-depth-gamma-delta-lower-bound}
 \operatorname{trdeg}_{\Qbar}
 \Qbar\bigl(e,\gamma,\delta,\EinDepth{2}(1)\bigr)\geq3.
\end{equation}
\end{corollary}

\begin{proof}
The first triple is the case \(s=1\), \(\lambda_1=1\), and \(R=2\) of
Theorem~\ref{thm:ein-depth-multipoint}.  The exact identity
\[
 2\EinDepth{2}(1)+2T=\gamma^2+\frac{\pi^2}{6}
\]
turns the first triple into the second by replacing the third coordinate
with twice itself.  Finally
\(\eta=\gamma+\delta/e\), so the first triple is contained in the field
on the left of \eqref{eq:ein-depth-gamma-delta-lower-bound}.
No individual irrationality or transcendence conclusion for \(\gamma\)
or \(\delta\) is asserted.
\end{proof}

\begin{remark}[Classical lineage and scope]
The functions in \eqref{eq:ein-depth-definition} belong to Hardy's
classical polyexponential lineage.  In the notation
\[
 \mathfrak e_p(x,a)=\sum_{n\geq0}\frac{x^n}{n!(n+a)^p},
\]
one has
\(\EinDepth{r}(z)=z\mathfrak e_{r+1}(-z,1)\), whereas
\(\mathfrak e_1(x,1)=(e^x-1)/x\).  Thus \(a=1\) is a resonant integral
shift whose first coordinate is already elementary.  Shidlovskii proved an all-order
algebraic-independence theorem for a nearby shifted family with a
nonintegral rational shift; see \cite[Sec.~7.3.3]{Shidlovskii1989} and the
account \cite[Sec.~9]{Boyadzhiev2007}.  The argument above treats the
regularized integral-shift endpoint needed for \(\Ein\), simultaneously at
several algebraic dilation points, and identifies the complete toric
relation ideal.
It does not cover multiple-index polyexponentials, whose shuffle-type
relations lie outside the present one-level tower.  We make no broader
historical-priority claim.
\end{remark}
